\section{Sanciones}
\label{sec:sanciones}

Con el fin de resguardar el trabajo colaborativo y la responsabilidad individual, se establecen las siguientes sanciones internas para los integrantes que incurran en incumplimientos.

\subsection{Inasistencia a reuniones}

\begin{itemize}[label=--]
    \item Primera inasistencia sin aviso previo: llamado de atención formal por parte del líder, registrado por escrito en el canal oficial del grupo.
    \item Segunda inasistencia sin aviso previo: el integrante deberá asumir una carga adicional de trabajo, definida por el resto del grupo, como compensación.
    \item Tercera inasistencia sin aviso previo: se informará la situación al docente, dejando constancia de la reiteración del incumplimiento.
\end{itemize}

\subsection{Inasistencia a la presentación final}

\begin{itemize}[label=--]
    \item La inasistencia injustificada a la presentación o defensa oral del trabajo se considerará una falta grave.
    \item En este caso, el integrante ausente podrá ser reportado directamente al docente como no participante en dicha etapa del trabajo, sin perjuicio de la nota grupal obtenida por el resto del equipo.
    \item Si la inasistencia cuenta con justificación válida (certificado médico u otra causa acreditable) y se informa con anticipación razonable, no se aplicará sanción.
\end{itemize}

\subsection{No entrega de avances o incumplimiento de tareas asignadas}

\begin{itemize}[label=--]
    \item Si un integrante no entrega su parte del trabajo dentro del plazo interno acordado, sin aviso ni justificación, se le enviará un recordatorio formal con un plazo final de 24 horas.
    \item Si transcurrido ese plazo el trabajo sigue sin entregarse, el resto del grupo podrá redistribuir la tarea y dejar constancia escrita del incumplimiento (registro en el chat grupal y en el historial de GitHub).
    \item La reiteración de incumplimientos (dos o más tareas no entregadas) faculta al grupo a informar la situación al docente, solicitando que la evaluación considere la participación individual real de cada integrante.
\end{itemize}

\subsection{Trabajo de baja calidad o no realizado conforme a lo acordado}

\begin{itemize}[label=--]
    \item El trabajo entregado que no cumpla con los estándares mínimos acordados por el grupo deberá ser corregido por el propio integrante dentro de un plazo adicional definido por el líder.
    \item De persistir la situación, se aplicarán las mismas medidas indicadas en la sección anterior (redistribución de tareas y aviso al docente).
\end{itemize}
